\chapter{O que a dopamina diz de verdade}
% versão completa: chapters/07_dofaalt.tex

No capítulo anterior a dopamina era simples: um número, ``melhor ou pior do que o esperado''. É uma primeira aproximação, e ela funciona bem. Mas os experimentos dos últimos anos mostraram que por esse ``fio'' passa mais do que um número. Este capítulo trata de onde o quadro é mais complicado e de onde os cientistas ainda discutem.

\section*{Otimistas e pessimistas}

Os neurônios de dopamina não são todos iguais. Uns reagem mais às surpresas agradáveis do que às desagradáveis: são os otimistas. Outros, o contrário: são os pessimistas. Se cada neurônio aprende pela sua própria regra, o otimista vai aprender como será a recompensa no melhor caso, e o pessimista, no pior. Juntos eles registram não uma expectativa média, mas toda a dispersão dos resultados possíveis, como um apostador que conhece não só o favorito, mas as chances de cada um. A ideia concorda com os experimentos; para que o cérebro precisa da dispersão inteira, por enquanto, é hipótese.

\pencilfig{7}{\pencilart{7}}{Cada neurônio de dopamina guarda o seu ponto na curva das recompensas possíveis: juntos, eles lembram a curva inteira.}

\section*{Não só ``melhor'', mas também ``importante''}

Uma parte dos neurônios de dopamina também dispara com eventos desagradáveis, mas importantes: um som alto, um perigo. Parece que pela rede de dopamina passa não só o sinal da surpresa, mas também a sua relevância --- um sinal de ``preste atenção''.

\section*{Dois sinais num só fio}

As rajadas rápidas ensinam. Já o nível médio lento de dopamina, que muda em minutos, parece falar de outra coisa: de quão rico está o ambiente agora. Se há muita recompensa por perto, o tempo é caro --- e o animal age mais depressa. Os experimentos confirmam que o nível de dopamina está ligado ao vigor das ações. Um engenheiro chamaria isso de separação de sinais por frequência: as altas frequências levam uma coisa, as baixas, outra.

\section*{A rampa perto do alvo}

Quando o animal se aproxima da recompensa, muitas vezes a dopamina não dá uma rajada, e sim sobe suavemente. Uma explicação: é a própria expectativa, o sinal ``o alvo está perto, acelere''. Outra mantém a ideia de surpresa: de longe o animal avalia a situação de forma borrada, e cada vez mais nítida conforme se aproxima, e cada refinamento é uma pequena surpresa agradável. Um experimento elegante num corredor virtual, em que o animal era transportado instantaneamente para mais perto do alvo, fala a favor da segunda: a dopamina respondeu com uma rajada, como convém a uma surpresa.

\section*{Olhando para trás}

Há também uma alternativa ousada: a dopamina responde não à pergunta ``o que vai acontecer?'', e sim à pergunta ``o que causou a recompensa?'' --- olha para o passado, não para o futuro. Os experimentos em que as duas teorias preveem coisas diferentes, por enquanto, se contradizem. É uma disputa científica viva.

\section*{Um vetor em vez de um número}

Por fim, a dopamina também reage à troca do \emph{sabor} da recompensa, quando o valor dela não mudou, e ajuda a aprender sem recompensa nenhuma. Parece que não é um erro só, e sim um feixe inteiro de erros --- um para cada traço do que aconteceu.

A conclusão do capítulo: a maioria das teorias novas não anula o quadro simples, e sim o amplia. A dopamina não é um fio só, mas um pequeno chicote de vários fios. Só o ``olhar para trás'' o contesta de verdade.

\fullversion{7}
